\documentclass[a4paper,12pt]{article}
\usepackage[utf8]{inputenc}
\usepackage[T1]{fontenc}
\usepackage{graphicx}
\usepackage{amsmath}
\usepackage{geometry}
\usepackage{booktabs}
\usepackage{hyperref}
\usepackage{caption}

\geometry{margin=1in}

\title{W2 Baseline Report: Active Learning for Medical Image Annotation}
\author{Antigravity AI Assistant}
\date{May 15, 2026}

\begin{document}

\maketitle

\section{Dataset Analysis}
The report is based on the \textbf{PneumoniaMNIST} dataset (MedMNIST), which contains pediatric chest X-ray images for binary pneumonia classification.

\subsection{Descriptive Statistics}
As documented in the \texttt{RL\_ActiveLearning\_SOTA.ipynb} data pipeline:
\begin{itemize}
    \item \textbf{Total Samples:} 5,856 images.
    \item \textbf{Splits:} 
    \begin{itemize}
        \item Training Set: 4,708 samples.
        \item Validation Set: 524 samples.
        \item Test Set: 624 samples.
    \end{itemize}
    \item \textbf{Image Dimensions:} $28 \times 28$ pixels, single channel (grayscale).
    \item \textbf{Missing Data:} Analysis confirms zero missing pixels/labels in all splits.
\end{itemize}

\subsection{Class Imbalance}
The training set exhibits a significant imbalance, which is a critical factor for model training and Active Learning strategy design:
\begin{itemize}
    \item \textbf{Class 0 (Normal):} 1,214 samples (25.8\%).
    \item \textbf{Class 1 (Pneumonia):} 3,494 samples (74.2\%).
    \item \textbf{Imbalance Ratio:} $\approx 2.87\times$.
\end{itemize}

\subsection{Preprocessing Requirements}
The pipeline implements the following transformations:
\begin{enumerate}
    \item \textbf{ToTensor:} Conversion of raw pixel values to floating-point tensors.
    \item \textbf{Normalization:} Application of \texttt{mean=[0.5]} and \texttt{std=[0.5]} to standardize intensities within the range $[-1, 1]$.
\end{enumerate}

\begin{figure}[h]
    \centering
    \fbox{Sample Visualizations: pneumoniamnist\_samples.png}
    \caption{Samples of normal vs. pneumonia X-ray images (extracted via \texttt{plot\_samples}).}
\end{figure}

\newpage

\section{Exploratory Data Analysis (EDA)}
The EDA focused on class distributions and visual clarity at low resolutions.

\subsection{Visual Inspection}
Visualized samples indicate that while the resolution is low ($28 \times 28$), lung opacities associated with pneumonia are still distinguishable as higher-intensity regions (white cloudy areas) compared to the clear lung fields in normal samples.

\subsection{Distribution Insights}
The class distribution analysis (performed via \texttt{plot\_class\_distribution}) highlights the dominance of the pneumonia class. This suggests that a naive model might achieve high accuracy by over-predicting the majority class, making metrics like AUC and Recall more representative of true clinical performance.

\begin{figure}[h]
    \centering
    \fbox{Train Class Distribution: class\_dist.png}
    \caption{Class frequency plot showing the 74.2\% majority pneumonia class.}
\end{figure}

\section{Baseline Description}
The baseline model is a custom Convolutional Neural Network (\textbf{MedicalCNN}) optimized for training on limited labeled data.

\subsection{Architecture Details}
The \texttt{MedicalCNN} architecture consists of:
\begin{itemize}
    \item \textbf{3 Convolutional Blocks:} Filters of size 16, 32, and 64 respectively, each followed by Batch Normalization and ReLU activation.
    \item \textbf{Downsampling:} Max Pooling ($2\times2$) and a final Adaptive Average Pooling ($4\times4$).
    \item \textbf{Classifier:} Two Fully Connected layers (128 units and 2 units) with interleaved Dropout (0.3 and 0.15) to mitigate overfitting.
\end{itemize}
Total trainable parameters: \textbf{154,978}.

\subsection{Training Configuration}
\begin{itemize}
    \item \textbf{Loss:} Cross-Entropy Loss.
    \item \textbf{Optimizer:} Adam with weight decay ($1\times 10^{-4}$).
    \item \textbf{Scheduler:} Cosine Annealing learning rate.
\end{itemize}

\newpage

\section{Baseline Metrics and Interpretation}
The baseline results were computed by training the \texttt{MedicalCNN} on an initial random seed of \textbf{100 labeled samples}.

\subsection{Performance Results}
The following metrics were recorded for the initial baseline:

\begin{table}[h]
\centering
\begin{tabular}{lc}
\toprule
\textbf{Metric} & \textbf{Baseline Value (n=100)} \\
\midrule
Validation AUC & 0.8752 \\
Test AUC & 0.8826 \\
Test Recall & 1.0000 \\
\bottomrule
\end{tabular}
\caption{Baseline model metrics before Active Learning iterations.}
\end{table}

\subsection{Results Interpretation}
\begin{itemize}
    \item \textbf{Strong Generalization:} Despite being trained on only 100 samples (2.1\% of the training pool), the model achieves a Test AUC of 0.8826, indicating a high ability to rank pneumonia cases above normal cases.
    \item \textbf{Clinical Safety:} The Recall of 1.00 is particularly noteworthy. It indicates that the model did not miss a single case of pneumonia in the test set at this threshold, which is critical for screening applications.
    \item \textbf{Training Stability:} The validation AUC (0.8752) is consistent with the test performance, suggesting the regularization (Dropout and BatchNorm) is effective.
\end{itemize}

\subsection{Discussion}
The baseline performance establishes a high bar for subsequent Active Learning strategies. While the initial recall is perfect, the AUC suggests room for improvement in reducing False Positives (Normal cases predicted as Pneumonia). The significant class imbalance likely contributes to the model's sensitivity towards Pneumonia. The SOTA RL-based approach implemented in the notebook aims to further optimize these metrics by selecting the most informative samples beyond the initial 100-sample seed.

\end{document}
